\begin{problem}[第34届物理竞赛复赛][弗兰克-赫兹实验]{77}
    1914 年，弗兰克-赫兹用电子碰撞原子的方法使原子从低能级激发到高能级，从而证明了原子能级的存在。加速电子碰撞自由的氢原子，使某氢原子从基态激发到激发态。该氢原子仅能发出一条可见光波长范围（$400\,\mathrm{nm} \sim 760\,\mathrm{nm}$）内的光谱线。仅考虑一维正碰。
    \begin{subproblem}
        求该氢原子能发出的可见光的波长；
    \end{subproblem}
    \begin{subproblem}
        求加速后电子动能 $E_{k}$ 的范围；
    \end{subproblem}
    \begin{subproblem}
        如果将电子改为质子，求加速质子的加速电压的范围。
    \end{subproblem}
    已知 $hc = 1240\,\mathrm{nm \cdot eV}$，其中 $h$ 为普朗克常量，$c$ 为真空中的光速；质子质量近似为电子质量的 $1836$ 倍，氢原子在碰撞前的速度可忽略。
\end{problem}